% FORMALIZACION_REUNIDA de la proyección del calendario y la representación
% de doce posiciones ya desarrolladas en la arquitectura TPK y el enlace Witt.
% Formalización explícita del cotejo de Gram y del transporte de marco.
\subsubsection{Projection du calendrier et représentation de ses douze positions}
\label{ii:dodeca:calendario}
Le calendrier de la section~\ref{ii:tpk-arquitectura} fournit une
coordonnée observable de l’état complet. Avec ses notations, nous écrivons
\(r_0(\theta)=r(\theta)-1\in\{0,\ldots,8\}\) et
\(b=\beta(\theta)\in\mathbb Z/12\mathbb Z\). La projection d’horloge est
\[
 Q_{\rm clk}(x)=\bigl(\beta(\theta(x)),r_0(\theta(x))\bigr).
\]
Pour un chemin admissible de \(n\) mises à jour, le transport induit est
\begin{equation}
 k(r_0,n)=\left\lfloor\frac{r_0+n}{9}\right\rfloor,
 \qquad
 T_n(b,r_0)=\bigl(b+k(r_0,n)\bmod12,\ (r_0+n)\bmod9\bigr).
 \label{ii:dodeca:transporte-reloj}
\end{equation}
Ici, \(n\in\mathbb N_0\) compte les mises à jour effectives ; le quotient
entier \(k\) enregistre les retours nonadiques traversés.

\begin{proposition}[Compatibilité du calendrier avec la composition des chemins]
\label{ii:dodeca:composicion}
La projection \(Q_{\rm clk}\), avec la longueur de chaque chemin, définit
un foncteur de la catégorie des chemins admissibles TPK vers la catégorie d’action
du calendrier~\eqref{ii:dodeca:transporte-reloj}. On a
\(T_{m+n}=T_mT_n\), \(T_0=I\),
\(T_9(b,r_0)=(b+1,r_0)\) et \(T_{108}(b,r_0)=(b,r_0)\).
\end{proposition}
\begin{proof}
En posant \(r'_0=(r_0+n)\bmod9\), la division euclidienne donne
\[
 k(r_0,n+m)=k(r_0,n)+k(r'_0,m).
\]
Cette identité conserve simultanément le résidu et le quotient lors de la concaténation
des chemins. La mise à jour de \(\theta\) établie dans la
section~\ref{ii:tpk-arquitectura} induit exactement
\eqref{ii:dodeca:transporte-reloj} ; les identités restantes s’obtiennent
par substitution. La catégorie d’arrivée a les paires \((b,r_0)\) pour
objets et les longueurs admissibles pour flèches, avec la somme comme composition.
\end{proof}

L’indice \(p=b+1\in\{1,\ldots,12\}\), en prenant le représentant
\(0\le b<12\), indique la position dodécaphasique. Un tour de
\(108\) mises à jour parcourt ses douze positions. La périodicité
démontrée correspond à cette projection : le chemin, la retenue accumulée et la
mémoire demeurent dans l’état source et poursuivent leur mise à jour.

\paragraph{Réalisation d’incidence à partir du projecteur angulaire.}
Les représentants \(r_pC_5\) et la représentation \(\rho\) de la
section~\ref{sec:ii-enlace-angular-witt} fixent la coordinatisation des douze
positions. Le projecteur déjà défini par
\eqref{eq:ii-witt-projector-character} produit
\begin{equation}
 \begin{gathered}
 \mathcal H_3=P_3^{\mathrm{ico}}\mathbb R^{12},\qquad v_p=2P_3^{\mathrm{ico}}e_p,
 \\[-2pt]
 (P_3^{\mathrm{ico}})_{pq}=\frac1{20}
 \sum_{\substack{g\in A_5\\gr_qC_5=r_pC_5}}\chi_3(g^{-1}).
 \end{gathered}
 \label{ii:dodeca:vertices-proyector}
\end{equation}
La dernière somme contient cinq termes par entrée et utilise les
représentants et les deux classes de cycles d’ordre cinq qui y sont fixés.

\begin{proposition}[Réalisation icosaédrique des positions dodécaphasiques]
\label{ii:dodeca:realizacion}
Les vecteurs \(v_p\) sont douze vecteurs unitaires distincts de
\(\mathcal H_3\). Leur incidence est
\(p\sim q\Longleftrightarrow\langle v_p,v_q\rangle=1/\sqrt5\).
L’application \(p\mapsto v_p\) est \(A_5\)-équivariante et son image
est isométrique à \(\Omega_{12}^{\rm ico}\).
\end{proposition}
\begin{proof}
La transitivité et \(\operatorname{Tr}P_3^{\mathrm{ico}}=3\) donnent
\((P_3^{\mathrm{ico}})_{pp}=1/4\) ; par idempotence,
\(\langle v_p,v_q\rangle=4(P_3^{\mathrm{ico}})_{pq}\) et \(\|v_p\|=1\).
L’évaluation des cinq termes de
\eqref{ii:dodeca:vertices-proyector} donne les paires antipodales
\[
 (1,4),\ (2,3),\ (5,8),\ (6,7),\ (9,12),\ (10,11).
\]
Pour les représentants \((1,2,5,6,9,10)\), leur matrice de Gram est
\begin{equation}
 I_6+\frac D{\sqrt5},\qquad
 D=\begin{pmatrix}
 0&-1&-1&1&-1&1\\
 -1&0&1&1&-1&-1\\
 -1&1&0&1&1&1\\
 1&1&1&0&-1&1\\
 -1&-1&1&-1&0&1\\
 1&-1&1&1&1&0
 \end{pmatrix},\qquad D^2=5I_6.
 \label{ii:dodeca:gram}
\end{equation}
Les autres produits scalaires s’obtiennent par antipodie.
Ainsi, hormis un sommet et son opposé, les produits sont
\(\pm1/\sqrt5\), et chaque sommet a cinq voisins. Le tableau suivant
explicite la comparaison avec la coordinatisation signée à six positions
\(\mathcal U=(\infty,0,1,2,3,4)\) :
\begin{equation}
\begin{array}{c|cc@{\qquad}c|cc}
p&u(p)&\sigma(p)&p&u(p)&\sigma(p)\\ \hline
1&\infty&+1&7&3&-1\\
2&0&-1&8&1&+1\\
3&0&+1&9&2&-1\\
4&\infty&-1&10&4&+1\\
5&1&-1&11&4&-1\\
6&3&+1&12&2&+1
\end{array}
\label{ii:dodeca:carta-firmada}
\end{equation}
\begin{sloppypar}
En effet, pour \(C=\operatorname{bal}(A_W)\), avec \(A_W\) donné dans
\eqref{exc:aw}, les entrées de \(D\) sont
\(D_{ab}=\sigma(p_a)\sigma(p_b)C_{u(p_a),u(p_b)}\).
Par conséquent, le tableau et~\eqref{ii:dodeca:gram} identifient
isométriquement \(v_p\) à
\(\sigma(p)\sqrt2 E_Ce_{u(p)}\), où
\(E_C=(I_6+C/\sqrt5)/2\). Les deux systèmes engendrent des espaces de dimension
trois et ont une matrice de Gram identique, ce qui prouve que l’application
linéaire induite est bien définie et isométrique. La réalisation
explicite de~\eqref{ii:ico:proyector-seis}, démontrée ci-après,
l’identifie à~\(\Omega_{12}^{\rm ico}\) et à son incidence.
Enfin, \(P_3^{\mathrm{ico}}\rho(g)=\rho(g)P_3^{\mathrm{ico}}\) implique
\(v_{g\cdot p}=\rho(g)v_p\). La bijection des douze positions donne
\(\operatorname{Stab}_{A_5}(v_p)=r_pC_5r_p^{-1}\).
\end{sloppypar}
\end{proof}

On obtient ainsi la composition représentationnelle
\begin{equation}
 x\longmapsto Q_{\rm clk}(x)\longmapsto
 p=\beta(\theta(x))+1\longmapsto v_p.
 \label{ii:dodeca:composicion-representacional}
\end{equation}
L’expression signée de cette composition est
\(r_\pm^{\rm rep}(x)=(u(p),\sigma(p))\), déterminée par le
tableau~\eqref{ii:dodeca:carta-firmada}. Le tableau est un choix explicite
d’alignement entre coordinatisations ; les actions de \(A_5\) fournissent leurs
alignements équivalents. Le registre complet est conservé dans le
domaine de~\eqref{ii:dodeca:composicion-representacional}.

\paragraph{Lien avec le porteur tensoriel à cinq composantes.}
Dans \(\mathcal H_3\), nous formons
\(Q_p=v_pv_p^*-I_{\mathcal H_3}/3\). Soit \(J=\mathbf1\mathbf1^*\)
et soit \(\mathsf A\) la matrice de la permutation antipodale précédente.
Le produit de Frobenius et~\eqref{ii:dodeca:gram} donnent
\begin{equation}
 \begin{aligned}
 \operatorname{Gram}(Q_p)
 &=(4P_3^{\mathrm{ico}})\circ(4P_3^{\mathrm{ico}})-\frac13J
 =\frac45(I+\mathsf A)-\frac2{15}J,
 \\
 P^{\rm ico}_5&=\frac58\operatorname{Gram}(Q_p)
 =\frac12(I+\mathsf A)-\frac1{12}J.
 \end{aligned}
 \label{ii:dodeca:puente-tensorial}
\end{equation}
Le symbole \(\circ\) désigne le produit entrée par entrée. La dernière
expression est le projecteur orthogonal sur les coordonnées paires sous
antipodie et orthogonales à \(\mathbf1\), de dimension cinq. Son caractère
est \((5,1,-1,0,0)\) : sur les six paires antipodales, les nombres de
points fixes sont \((6,2,0,1,1)\), et l’on retire la composante constante.
Il coïncide donc avec le projecteur central de
\eqref{ii:pent:eq:gravity-p5}. Cette identité relie la construction
tridimensionnelle à \(\Gamma_0\) et \(\Gamma_5\), définis dans
\eqref{ii:pent:eq:gravity-Gamma0} et~\eqref{ii:pent:eq:gravity-gamma5},
au moyen de l’opération quadratique explicite \(v_p\mapsto Q_p\).

\subsubsection{Transport du repère représentationnel}
\label{ii:dodeca:marco}
Soit \(s=(1\,2\,\cdots\,12)\) et soit \(S_ce_p=e_{s(p)}\).
Un retour nonadique du calendrier envoie \(p\) sur \(s(p)\). Le transport
conjoint du repère, de la représentation et du projecteur s’exprime par
\begin{equation}
 \begin{gathered}
 P_{3,k}^{\mathrm{ico}}=S_c^kP_3^{\mathrm{ico}}S_c^{-k},\qquad
 \rho^{(k)}(g)=S_c^k\rho(g)S_c^{-k},\\
 v^{(k)}_{s^k(p)}=2P_{3,k}^{\mathrm{ico}}e_{s^k(p)}=S_c^kv_p.
 \end{gathered}
 \label{ii:dodeca:marco-movil}
\end{equation}
La dernière égalité se prouve par substitution ; l’orthogonalité de
\(S_c\) conserve tous les produits scalaires et transporte
l’incidence. Pour un chemin de longueur \(n\), on prend
\(k=k(r_0,n)\). L’identité de retenue de la
proposition~\ref{ii:dodeca:composicion} prouve la compatibilité de
ces transports avec la concaténation. Bien que \(S_c^{12}=I\),
les douze tours mettent à jour la mémoire du chemin source.

L’égalité~\eqref{ii:dodeca:marco-movil} établit la covariance entre
repères. L’action active sur une incidence fixe a un contenu
supplémentaire : elle nécessite d’identifier le transport de l’état complet et son
action dans cette coordinatisation. En particulier, \(s\) est d’ordre douze, tandis
que les ordres des éléments de \(A_5\) sont \(1,2,3,5\).
De même, le stabilisateur d’un sommet est d’ordre cinq, et celui de sa
paire antipodale est d’ordre dix ; l’involution de la coordinatisation signée
échange les deux sommets. Ces distinctions conservent le sens
de l’orientation, du feuillet et du retour dans les propriétés (i)--(iv) : la
réalisation représentationnelle démontrée et la descente du chemin
exceptionnel complet se relient au moyen de ces propriétés, avec leurs
domaines et fibres déclarés.
